The idea of representing RC operations as explicit instructions so as to
optimize them via static analysis is described as early as
\citet{Barth:1977:SGC:359636.359713}.
\citet{Schulte:1994} describes a system with many
features similar to ours. %% , as we will detail below.
In general, Schulte's language is much simpler than ours, with a single list type
as the only non-primitive type, and no higher-order functions. He does
not give a formal dynamic semantics % or static type system
for his system. He gives an algorithm for inserting RC instructions that, like ours,
has an on-the-fly optimization for omitting \kw{inc} instructions if a
variable is already dead and would immediately be decremented
afterwards. Schulte briefly discusses how RC operations can be
minimized by treating some parameters as ``nondestructive'' in the
sense of our borrowed references. In contrast to our inference of
borrow annotations, Schulte proposes to create one copy of a function
for each possible destructive/nondestructive combination of parameters
(i.e.\ exponential in the number of (non-primitive) parameters) and
to select an appropriate version for each call site of the function. Our
approach never duplicates code.

Introducing destructive updates into pure programs has traditionally focused on
primitive operations like array updates~\cite{Hudak:1985:AUP:318593.318660}, particularly in the functional array
languages \textsc{Sisal}~\cite{sisal} and \textsc{SaC}~\cite{sac}. \citet{sac-reuse} propose an \verb!alloc_or_reuse!
instruction for \textsc{SaC} that can select one of multiple array candidates for reuse, but do not describe
heuristics for when to use the instruction.
\citet{Ferey:2016:CGU:2963372.2963387} describe how functional update operations
explicit in the source language can be turned into destructive updates using the
reference counter. In contrast, \citet{Schulte:1994} presents a ``reusage'' optimization that has an effect similar to the one
obtained with our \kw{reset}/\kw{reuse} instructions. In particular, it is
independent of a specific surface-level update syntax.
However, his optimization (transformation $T14$) is more restrictive
and is only applicable to a branch of a $\kw{case}\ x$ if $x$ is dead at the beginning of the branch.
% Although Schulte demonstrates his optimization can produce code for the list append function where
% destructive updates are performed when the list is not shared,
His optimization cannot handle the simple $\textit{swap}$ described earlier,
let alone more complex functions such as the red black
tree re-balancing function $\textit{balance}_1$.

While not a purely functional language, the Swift programming
language\footnote{\url{https://developer.apple.com/swift/}} has directly
influenced many parts of our work. To the best of our knowledge, Swift was the
first non-research language to use an intermediate representation with explicit RC
instructions, as well as the idea of (safely) avoiding RC operations via ``borrowed''
parameters (which are called ``+0'' or ``guaranteed'' in Swift), in its
implementation. While Swift's primitives may also elide copies when given a
unique reference, no speculative destructive updates are introduced for
user-defined types, but this may not be as important for an impure language as
it is for Lean. Parameters default to borrowed in Swift, but the compiler may
locally change the calling convention inside individual modules.

\citet{Baker:1994:MRC:185009.185016} describes optimizing reference counting by
use of \emph{two} pointer kinds, a standard one and a \emph{deferred increment}
pointer kind. The latter kind can be copied freely without adding RC operations,
but must be converted into the standard kind by incrementing it before storing it in an object or returning it. The two kinds
are distinguished at runtime by pointer tagging. Our borrowed references
can be viewed as a static refinement of this idea. Baker then describes an
extended version of deferred-increment he calls \emph{anchored} pointers that
store the stack level (i.e.\ the lifetime) of the standard pointer they have been
created from. Anchored pointers do not have to be converted to the standard kind if returned
from a stack frame above this level. In order to statically approximate this
extended system, we would need to extend our type system
with support for some kind of \emph{lifetime annotations} on return types
as featured in Cyclone~\cite{jim2002cyclone} and
Rust~\cite{Matsakis:2014:RL:2663171.2663188}.

%% In our notation, his optimization could be presented as transforming
%% \[ \Case x F_1 (\Dec x; C[\Ctor_2 ])\]

%% \citet{chirimar_gunter_riecke_1996} give both a natural and a reference-counting
%% big-step semantics on a linearly typed language and prove their equivalence.
%% Their linearly typed language is not a simplified intermediate representation but the
%% source language. We have discussed how their proofs could be adapted to our extended type system and
%% semantics in \cref{sec:ty}. \citet{Ferey:2016:CGU:2963372.2963387} give a
%% natural and a RC big-step semantics on a non-linear intermediate representation
%% without explicit RC operations, which are directly part of the RC semantics instead.
%% \citet{Hudak:1986:SMR:319838.319876} gives a denotational semantics of reference counting.

\citet{dynamic-atomicity} optimize Swift's reference counting scheme
by using a single bit to tag objects possibly shared between multiple
threads, much like our approach. However, because of mutability, every
single store operation must be intercepted to (recursively) tag objects
before becoming reachable from an already tagged object.
\citet{biasedRC} remove the need for tagging by extending every object header
with the ID of the thread $T$ that allocated the value, and two
reference counters: a shared one that requires atomic operations, and
another one that is only updated by $T$.
Thanks to immutability, we can make use of the simpler scheme without
introducing store barriers during normal code generation. Object tagging
instead only has to be done in threading primitives.

%%% Local Variables:
%%% mode: latex
%%% TeX-master: "refcount"
%%% End:
